\section*{第 \ref{ch:g12:continental-crust} 章——陆壳}

\begin{solution}{exo:g12:continental-crust:1}
陆壳：厚度 30 至 \qty{70}{km}，平均约 \qty{41}{km}；
花岗质岩石、片麻岩及其他变质岩和\omterm{def:g11:igneous-rocks:volcanic-plutonic}{深成岩}，位于
沉积\omterm{def:g11:fossil-fuel-geology:reservoir}{盖层}之下；平均二氧化硅含量约 \qty{61}{\percent}；密度约
\qty{2830}{kg/m^3}；最老可达 40.3 亿年（颗粒为 44 亿年）。\omterm{def:g11:ocean-ridges:oceanic-lithosphere}{洋壳}：
厚约 \qty{7}{km}；玄武岩和\omterm{def:g11:ocean-ridges:gabbro}{辉长岩}；平均二氧化硅含量约 \qty{50}{\percent}；密度约
\qty{3000}{kg/m^3}；不足 2 亿年。
\end{solution}

\begin{solution}{exo:g12:continental-crust:2}
\qty{60.6}{\percent}：\omterm{def:g11:igneous-rocks:felsic-mafic}{中性}，类似\omterm{def:g11:igneous-rocks:rock-names}{闪长岩}。\qty{66.6}{\percent}：
\omterm{def:g11:igneous-rocks:felsic-mafic}{长英质}，类似花岗岩。
\end{solution}

\begin{solution}{exo:g12:continental-crust:3}
约 \qty{3000}{kg/m^3} 的\omterm{def:g11:ocean-ridges:oceanic-lithosphere}{洋壳}位于\omterm{def:g11:lithosphere-asthenosphere:lithosphere}{岩石圈}地幔之上，
而古老、寒冷的\omterm{def:g11:ocean-ridges:oceanic-lithosphere}{大洋岩石圈}整体足够致密，能够下沉。
约 \qty{2830}{kg/m^3} 的陆壳远比地幔轻
（约 \qty{3380}{kg/m^3}），浮力过大：大陆抵达海沟时
便会阻止\omterm{def:g11:subduction-zones:subduction}{俯冲}。
\end{solution}

\begin{solution}{exo:g12:continental-crust:4}
（1）古老而稳定的大陆\omterm{def:g12:continental-crust:craton}{克拉通}，已不再变形。（2）超过
10 亿年。（3）变质岩和\omterm{def:g11:igneous-rocks:volcanic-plutonic}{深成岩}，是
已消失山脉的深部根基。
\end{solution}

\begin{solution}{exo:g12:continental-crust:5}
$4030/190 \approx 21$ 倍。洋底在洋脊处生成，并在
\omterm{def:g11:subduction-zones:subduction}{俯冲带}消亡，没有哪一段洋底能存在远超过
2 亿年。
\end{solution}

\begin{solution}{exo:g12:continental-crust:6}
（1）$0.306 \times 210.4 \approx \num{64.4}\text{ 百万 }\unit{km^2}$。
（2）$210.4 - 64.4 = \num{146.0}\text{ 百万 }\unit{km^2}$。
（3）$4\pi \times 6371^2 \approx \qty{5.10e8}{km^2}$；$210.4/510.1 \approx
0.41$：\qty{41}{\percent}。
\end{solution}

\begin{solution}{exo:g12:continental-crust:7}
（1）$41100 \times 2830 \approx \qty{1.16e8}{kg}$。（2）$7100 \times 3000 =
\qty{2.13e7}{kg}$。（3）约为 5.5 倍。
\end{solution}

\begin{solution}{exo:g12:continental-crust:8}
\omterm{def:g12:continental-crust:root}{山根}：$2.5 \times 2830/550 \approx \qty{12.9}{km}$。地壳增厚
$2.5 \times 3380/550 \approx \qty{15.4}{km}$：合计约 \qty{50}{km}，
位于山带之下所见的 50 至 \qty{70}{km} 范围之内。
\end{solution}

\begin{solution}{exo:g12:continental-crust:9}
（1）约 $0.103 + 0.100 + 0.102 \approx 0.30$：\qty{30}{\percent}。
（2）最年轻的年龄段，0 至 1000 万年（约 \qty{10}{\percent}），
几乎与 2000 至 3000 万年段并列。（3）老洋底已在\omterm{def:g11:subduction-zones:subduction}{俯冲带}
消亡：洋底越老，消亡的部分就越多。
\end{solution}

\begin{solution}{exo:g12:continental-crust:10}
$120 = 0.694\,h - 0.306 \times 100$，故 $h = 150.6/0.694 \approx
\qty{217}{m}$。仅统计陆地会漏掉低洼的、被海水淹没的大陆边缘，
使平均值明显偏高。
\end{solution}

\begin{solution}{exo:g12:continental-crust:11}
岛屿：熔岩之下为\omterm{def:g11:ocean-ridges:gabbro}{辉长岩}，\omterm{def:g11:seismic-structure:moho}{莫霍面}位于 \qty{8}{km}，属\omterm{def:g11:ocean-ridges:oceanic-lithosphere}{洋壳}
（岛屿的海拔不改变这一点）。平原：片麻岩，
\omterm{def:g11:seismic-structure:moho}{莫霍面}位于 \qty{38}{km}，属陆壳。
\end{solution}

\begin{solution}{exo:g12:continental-crust:12}
（1）精确值：$\tfrac{4}{3}\pi(6371^3 - 6329.9^3) \approx
\qty{2.0829e10}{km^3}$；近似值：$4\pi \times 6371^2 \times 41.1
\approx \qty{2.0964e10}{km^3}$。（2）近似值偏大约
\qty{0.65}{\percent}，接近 $H/R \approx 0.0065$：对于薄层而言
误差可以忽略不计。
\end{solution}

\begin{solution}{exo:g12:continental-crust:13}
（1）$60.6 = (66.6 + 2x)/3$，故 $x = (181.8 - 66.6)/2 = \qty{57.6}{\percent}$。
（2）\omterm{def:g11:igneous-rocks:felsic-mafic}{中性}。（3）深部陆壳的二氧化硅含量低于
上部陆壳，\omterm{def:g11:igneous-rocks:felsic-mafic}{镁铁质}矿物更多，富含铁、镁和钙。
\end{solution}

\begin{solution}{exo:g12:continental-crust:14}
（1）$3 \times 10^8/3 = 10^8$ 年，即 1 亿年。（2）与
洋底年龄的量级相同：约四分之三的洋底年龄小于 1 亿年，
没有超过约 1.9 亿年的。（3）其消亡速度必须与生成速度一样快：
它在\omterm{def:g11:subduction-zones:subduction}{俯冲带}重新没入地幔。
\end{solution}

\begin{solution}{exo:g12:continental-crust:15}
（1）$8.6 \times 10^9 / 4.0 \times 10^9 \approx \qty{2.2}{km^3}$ 每年。
（2）稳定速率只是一种可能：地壳可能在早期快速增长，
此后被循环再造，也可能是较晚才增长的。最古老的岩石和颗粒证明
早期已存在一些地壳，但无法说明其总量，因为古老地壳
可能被侵蚀或被带回地幔，留下的痕迹很少。
\end{solution}

\begin{solution}{pb:g12:continental-crust:1}
\textbf{1.} $4\pi \times 6371^2 \approx \qty{5.10e8}{km^2}$。

\textbf{2.} $210.4/510.1 \approx 0.41$：\qty{41}{\percent}。

\textbf{3.} $0.306 \times 210.4 \approx \num{64.4}\text{ 百万 }\unit{km^2}$。

\textbf{4.} $510.1 - 210.4 \approx \num{299.7}\text{ 百万 }\unit{km^2}$，
\qty{59}{\percent}。

\textbf{5.} 大陆边缘位于浅海之下，
即位于大陆架上。

\textbf{6.} 厚度 \qty{41}{km} 与
半径 \qty{6371}{km} 相比很小：曲率引起的体积变化小于
百分之一。

\textbf{7.} $210.4 \times 10^6 \times 41.1 \approx \qty{8.65e9}{km^3}$。

\textbf{8.} \qty{8.65e18}{m^3}。

\textbf{9.} $210.4 \times 10^6 \times 36 \approx \qty{7.57e9}{km^3}$。

\textbf{10.} $(8.65 - 7.57)/7.57 \approx 0.14$：\qty{14}{\percent}。

\textbf{11.} $8.65 \times 10^{18} \times 2830 \approx \qty{2.45e22}{kg}$。

\textbf{12.} $41100 \times 2830 \approx \qty{1.16e8}{kg}$。

\textbf{13.} $299.7 \times 10^6 \times 7.1 \approx \qty{2.13e9}{km^3}$，
即 \qty{2.13e18}{m^3}；乘以 3000，约为 \qty{6.4e21}{kg}。

\textbf{14.} \omterm{def:g12:continental-crust:continental-crust}{大陆地壳}约重 3.8 倍：
它的厚度几乎为六倍，足以抵消面积和密度较小的影响。

\textbf{15.} \omterm{def:g12:continental-crust:continental-crust}{大陆地壳}由花岗质岩石组成，富含石英和长石，它们是
造岩矿物中最轻的一类；大洋地壳由玄武岩和
\omterm{def:g11:ocean-ridges:gabbro}{辉长岩}组成，富含辉石和橄榄石。

\textbf{16.} $2.45 \times 10^{22}/5.97 \times 10^{24} \approx 0.0041$。

\textbf{17.} $3 \times 0.41 \times (41.1/6371) \times (2830/5513) \approx
0.0041$：结果相同。

\textbf{18.} $6.4 \times 10^{21}/5.97 \times 10^{24} \approx 0.0011$：约为
千分之一。

\textbf{19.} 约 \qty{0.4}{\percent}。

\textbf{20.} \omterm{def:g12:continental-crust:continental-crust}{大陆地壳}约占\textbf{\qty{0.4}{\percent}
的地球质量}。
\end{solution}
